% Source: https://artofproblemsolving.com/community/c3h1349348_geometry
% Stable identity: aops-topic-1349348-q01
Let the rhombus have side $x$, so $AF=FC=x$ and $FB=16-x$.

\begin{asy}
usepackage("amsmath");
size(6.5cm);
defaultpen(fontsize(10pt));
pair A=(0,0),B=(16,0),C=(16,12),D=(0,12),F=(12.5,0),E=(3.5,12);
draw(A--B--C--D--cycle,linewidth(0.9));
draw(A--F--C--E--cycle,blue+linewidth(0.9));
draw(A--C,gray+dashed); draw(E--F,red);
dot(E);dot(F);
label("$A$",A,SW); label("$B$",B,SE); label("$C$",C,NE); label("$D$",D,NW);
label("$F$",F,S); label("$E$",E,N);
label("$x$",(A+F)/2,S,blue); label("$16-x$",(F+B)/2+(0.4,-1.3),S); label("$12$",(B+C)/2,E);
label("$x$",(F+C)/2,W,blue);
label("\mbox{}",(min(currentpicture,user=true)+max(currentpicture,user=true))/2,black); // keeps the text after the figure black when LaTeX typesets the labels inline
\end{asy}

In the right triangle $FBC$,
\[
(16-x)^2+12^2=x^2\iff 400-32x=0\iff x=\dfrac{25}{2}.
\]
The diagonal $AC$ of the rectangle has length $\sqrt{16^2+12^2}=20$. The area of the rhombus is base times height, $AF\cdot BC=\dfrac{25}{2}\cdot12=150$, and also half the product of its diagonals, $\dfrac{AC\cdot EF}{2}=10\,EF$. So $EF=\boxed{15}$.

(Coordinates confirm this: with $A=(0,0)$ and $C=(16,12)$ we get $F=\left(\dfrac{25}{2},0\right)$ and $E=\left(\dfrac{7}{2},12\right)$, and $EF=\sqrt{9^2+12^2}=15$.)
